% Generated draft: requires mathematical and source review.
\documentclass[11pt]{article}
\usepackage[margin=25mm]{geometry}
\usepackage{fontspec}
\setmainfont{Latin Modern Roman}
\usepackage{amsmath,amssymb,mathtools,unicode-math}
\setmathfont{Latin Modern Math}
\usepackage{enumitem,needspace,xurl,hyperref}
\setlength{\parindent}{0pt}
\setlength{\parskip}{5pt}
\setlength{\emergencystretch}{3em}
\newcommand{\N}{\mathbb{N}}
\newcommand{\Z}{\mathbb{Z}}
\newcommand{\Q}{\mathbb{Q}}
\newcommand{\R}{\mathbb{R}}
\newcommand{\C}{\mathbb{C}}
\newcommand{\F}{\mathbb{F}}
\newcommand{\A}{\mathbb{A}}
\newcommand{\PP}{\mathbb{P}}
\newcommand{\E}{\mathbb{E}}
\newcommand{\eps}{\varepsilon}
\newcommand{\dd}{\,\mathrm{d}}
\begin{document}
\section*{UnsolvedMath ID: 701557}
\textbf{Source problem code:} catalogue entry
\subsection{Definitions and mathematical statement}
For a group \(G\), write \({}^{g}h=ghg^{-1}\) for conjugation. The non-abelian tensor square \(G\otimes G\) is the group generated by symbols \(g\otimes h\), for \(g,h\in G\), subject to the relations
\[
(gg')\otimes h=({}^{g}g'\otimes{}^{g}h)(g\otimes h),
\qquad
g\otimes(hh')=(g\otimes h)({}^{h}g\otimes{}^{h}h')
\]
for all \(g,g',h,h'\in G\). Both actions used here are conjugation actions.

A group \(H\) is called linear if there exist a field \(K\), an integer \(d\geq1\), and an injective homomorphism \(H\hookrightarrow\operatorname{GL}_{d}(K)\). The source specifies neither a characteristic nor a common coefficient field for the groups under consideration. The questions below use this usual existential meaning of linearity; any intended requirement to preserve a particular field is unresolved by the supplied record.

For an integer \(n\geq2\), the braid group \(B_n\) is the group with presentation
\[
B_n=\left\langle \sigma_1,\ldots,\sigma_{n-1}\ \middle|\ 
\begin{aligned}
\sigma_i\sigma_j&=\sigma_j\sigma_i &&(1\leq i,j\leq n-1,\ |i-j|\geq2),\\
\sigma_i\sigma_{i+1}\sigma_i&=\sigma_{i+1}\sigma_i\sigma_{i+1} &&(1\leq i\leq n-2)
\end{aligned}
\right\rangle.
\]

\begin{enumerate}
\item[(a)] Is it true that, for every finitely generated linear group \(G\), the non-abelian tensor square \(G\otimes G\) is linear? Explicitly, does the existence of a field \(K\), an integer \(d\geq1\), and an embedding \(G\hookrightarrow\operatorname{GL}_{d}(K)\) imply the existence of a field \(L\), an integer \(e\geq1\), and an embedding \(G\otimes G\hookrightarrow\operatorname{GL}_{e}(L)\)?
\item[(b)] In particular, for every integer \(n>3\), is \(B_n\otimes B_n\) linear?
\end{enumerate}

The supplied source reports that \(B_3\otimes B_3\) is linear and that there exists a countable group \(H\) for which \(H\otimes H\) is not linear. In the latter assertion, the record does not assert that \(H\) is finitely generated or linear. These are preserved source-reported background assertions, not independently verified results. The registry labels the problem open; that label is supplied curation rather than an independently verified assessment.
\subsection{Short English statement}
The problem asks whether the non-abelian tensor square of every finitely generated linear group is linear, and asks specifically about braid groups \(B_n\) for integers \(n>3\). The source reports a positive result for \(B_3\) and a countable example with non-linear tensor square, without asserting the additional hypotheses of part (a) for that example. Coefficient fields are unspecified.
\subsection{Sources}
\begin{itemize}
\item Local UnsolvedMath registry, supplied complete source record: permanent ID \(701557\), exact source problem code \texttt{catalogue entry}, title \emph{Kourovka Notebook Problem 19.9}. The record labels the problem \texttt{open} and gives its update timestamp as \texttt{2026-10-16T14:00:00Z}; these are supplied registry metadata, not independent mathematical verification.
\item \emph{The Kourovka Notebook (Unsolved Problems in Group Theory)}, original problem 19.9, proposed by V. G. Bardakov and M. V. Neshchadim, 2018. The supplied background identifies the 19th issue (2018) and the 21st issue edition (2026), arXiv:1401.0300v48. The supplied collection description names E. I. Khukhro and V. D. Mazurov as editors. These edition and attribution details have not been independently checked.
\item R. Brown and J.-L. Loday, \emph{Topology} \textbf{26}, no. 3 (1987), 311--335. The source cites this reference for the definition of the non-abelian tensor square; no article title or URL is supplied.
\end{itemize}
\end{document}
